\documentclass[11pt,a4paper]{article}

% -------------------------------------------------
% Packages
% -------------------------------------------------
\usepackage[margin=1in]{geometry}
\usepackage{amsmath,amssymb,amsfonts,mathtools}
\usepackage{bm}
\usepackage{physics}
\usepackage{enumitem}
\usepackage{hyperref}

% -------------------------------------------------
% Commands
% -------------------------------------------------
\newcommand{\R}{\mathbb{R}}
\newcommand{\dd}{\mathrm{d}}
\newcommand{\PP}{\mathcal{P}}
\newcommand{\QQ}{\mathcal{Q}}

\begin{document}

\begin{center}
    {\Large\bf Poincaré Compactification of a Two-Dimensional
    Hindmarsh--Rose-Type Polynomial System}
\end{center}

\vspace{0.5cm}

We consider the two-dimensional Hindmarsh--Rose-type polynomial system
\begin{equation}
\boxed{
\begin{aligned}
\dot{x} &= y-a x^3+b x^2+I,\\
\dot{y} &= c-dx^2-y,
\end{aligned}}
\tag{HR}
\end{equation}
where $a,b,c,d,I$ are real parameters. Throughout the generic
cubic case, we assume
\begin{equation}
a\neq 0.
\end{equation}

The purpose of this section is to perform the complete Poincaré
compactification and to determine the singularities at infinity.

% =================================================
\section{1. Standard Form and Polynomial Degree}
% =================================================

Write the system in the standard form
\begin{equation}
\dot{x}=P(x,y),\qquad
\dot{y}=Q(x,y),
\end{equation}
where
\begin{equation}
\boxed{
P(x,y)=y-a x^3+b x^2+I
}
\end{equation}
and
\begin{equation}
\boxed{
Q(x,y)=c-dx^2-y.
}
\end{equation}

The degree of $P$ is
\begin{equation}
\deg P=3
\end{equation}
provided $a\neq0$, whereas
\begin{equation}
\deg Q=2
\end{equation}
provided $d\neq0$.

Hence the degree of the polynomial vector field is
\begin{equation}
\boxed{
n=\max\{\deg P,\deg Q\}=3.
}
\tag{1}
\end{equation}

Therefore, in the Poincaré compactification, the appropriate
time rescaling is
\begin{equation}
\boxed{
\frac{\dd}{\dd\tau}
=
v^{n-1}\frac{\dd}{\dd t}
=
v^2\frac{\dd}{\dd t}.
}
\tag{2}
\end{equation}

Equivalently,
\begin{equation}
\boxed{
\frac{\dd t}{\dd\tau}=v^2.
}
\end{equation}

% =================================================
\section{2. Homogeneous Components}
% =================================================

We decompose $P$ and $Q$ into homogeneous components.

For $P$,
\begin{equation}
P(x,y)
=
P_3(x,y)+P_2(x,y)+P_1(x,y)+P_0(x,y),
\end{equation}
where
\begin{align}
P_3(x,y)&=-a x^3,\\
P_2(x,y)&=b x^2,\\
P_1(x,y)&=y,\\
P_0(x,y)&=I.
\end{align}

Thus
\begin{equation}
\boxed{
P_3(x,y)=-a x^3.
}
\tag{3}
\end{equation}

For $Q$,
\begin{equation}
Q(x,y)
=
Q_3(x,y)+Q_2(x,y)+Q_1(x,y)+Q_0(x,y),
\end{equation}
where
\begin{align}
Q_3(x,y)&=0,\\
Q_2(x,y)&=-d x^2,\\
Q_1(x,y)&=-y,\\
Q_0(x,y)&=c.
\end{align}

Therefore
\begin{equation}
\boxed{
Q_3(x,y)=0.
}
\tag{4}
\end{equation}

Consequently, the highest-degree homogeneous part of the vector field is
\begin{equation}
\boxed{
(P_3,Q_3)
=
\left(-a x^3,0\right).
}
\tag{5}
\end{equation}

This highest-degree part determines the directions of the
singularities at infinity.

% =================================================
\section{3. Singular Directions at Infinity}
% =================================================

For a polynomial planar system of degree $n$, the singular
directions at infinity can be found by using the highest-degree
homogeneous components $P_n$ and $Q_n$.

In the $x$-chart, the singular directions satisfy
\begin{equation}
\boxed{
Q_n(1,u)-uP_n(1,u)=0.
}
\tag{6}
\end{equation}

For the present system, $n=3$, and therefore
\begin{equation}
P_3(1,u)=-a,
\qquad
Q_3(1,u)=0.
\end{equation}

Hence
\begin{align}
Q_3(1,u)-uP_3(1,u)
&=
0-u(-a)\\
&=au.
\end{align}

Thus
\begin{equation}
au=0.
\end{equation}

Under the generic assumption $a\neq0$,
\begin{equation}
\boxed{
u=0.
}
\tag{7}
\end{equation}

Therefore one singular direction at infinity is obtained in the
$x$-chart.

Similarly, in the $y$-chart, the singular directions satisfy
\begin{equation}
\boxed{
P_n(u,1)-uQ_n(u,1)=0.
}
\tag{8}
\end{equation}

Since
\begin{equation}
P_3(u,1)=-au^3,
\qquad
Q_3(u,1)=0,
\end{equation}
we obtain
\begin{align}
P_3(u,1)-uQ_3(u,1)
&=
-au^3.
\end{align}

Thus
\begin{equation}
-au^3=0,
\end{equation}
and hence
\begin{equation}
\boxed{
u=0.
}
\tag{9}
\end{equation}

Therefore the generic system has two singular directions at
infinity:
\begin{equation}
\boxed{
[1:0:0]
\qquad\text{and}\qquad
[0:1:0].
}
\tag{10}
\end{equation}

These correspond respectively to the $x$-axis and $y$-axis
directions at infinity.

% =================================================
\section{4. First Local Chart: the $x$-Chart}
% =================================================

Consider the local chart defined by
\begin{equation}
\boxed{
u=\frac{y}{x},
\qquad
v=\frac{1}{x}.
}
\tag{11}
\end{equation}

The inverse transformation is
\begin{equation}
\boxed{
x=\frac{1}{v},
\qquad
y=\frac{u}{v}.
}
\tag{12}
\end{equation}

The coordinate $v$ measures the approach to infinity:
\begin{equation}
|x|\to\infty
\quad\Longleftrightarrow\quad
v\to0.
\end{equation}

Thus the line
\begin{equation}
\boxed{
v=0
}
\tag{13}
\end{equation}
represents the set of directions at infinity in this chart.

The variable
\begin{equation}
u=\frac{y}{x}
\end{equation}
is the directional coordinate.

% =================================================
\subsection{4.1 Transformation of the $v$-Equation}
% =================================================

Since
\begin{equation}
x=\frac{1}{v},
\end{equation}
we have
\begin{equation}
\dot{x}
=
-\frac{\dot v}{v^2}.
\end{equation}

Since
\begin{equation}
\dot{x}=P(x,y),
\end{equation}
it follows that
\begin{equation}
-\frac{\dot v}{v^2}
=
P\left(\frac1v,\frac{u}{v}\right),
\end{equation}
and consequently
\begin{equation}
\boxed{
\dot v
=
-v^2P\left(\frac1v,\frac{u}{v}\right).
}
\tag{14}
\end{equation}

Now
\begin{align}
P\left(\frac1v,\frac{u}{v}\right)
&=
\frac{u}{v}
-a\frac{1}{v^3}
+b\frac{1}{v^2}
+I\\
&=
\frac{-a+bv+uv^2+Iv^3}{v^3}.
\end{align}

Therefore
\begin{align}
\dot v
&=
-v^2
\frac{-a+bv+uv^2+Iv^3}{v^3}\\
&=
\frac{a-bv-uv^2-Iv^3}{v}.
\end{align}

Thus
\begin{equation}
\boxed{
\dot v
=
\frac{a-bv-uv^2-Iv^3}{v}.
}
\tag{15}
\end{equation}

% =================================================
\subsection{4.2 Transformation of the $u$-Equation}
% =================================================

Since
\begin{equation}
u=\frac{y}{x},
\end{equation}
we have
\begin{equation}
\dot u
=
\frac{x\dot y-y\dot x}{x^2}.
\end{equation}

Using
\begin{equation}
x=\frac1v,\qquad
y=\frac{u}{v},
\end{equation}
we obtain
\begin{equation}
\boxed{
\dot u
=
v\left[
Q\left(\frac1v,\frac{u}{v}\right)
-
uP\left(\frac1v,\frac{u}{v}\right)
\right].
}
\tag{16}
\end{equation}

First,
\begin{align}
Q\left(\frac1v,\frac{u}{v}\right)
&=
c-d\frac1{v^2}-\frac{u}{v}\\
&=
\frac{cv^3-dv-uv^2}{v^3}.
\end{align}

Hence
\begin{align}
Q-uP
&=
\frac{cv^3-dv-uv^2}{v^3}\\
&\quad
-
u
\frac{-a+bv+uv^2+Iv^3}{v^3}\\
&=
\frac{
au-(d+bu)v-u(1+u)v^2+(c-Iu)v^3
}{v^3}.
\end{align}

Therefore,
\begin{equation}
\boxed{
\dot u
=
\frac{
au-(d+bu)v-u(1+u)v^2+(c-Iu)v^3
}{v^2}.
}
\tag{17}
\end{equation}

% =================================================
\subsection{4.3 Time Rescaling}
% =================================================

Since the degree of the system is $n=3$, we use
\begin{equation}
\boxed{
\frac{\dd}{\dd\tau}
=
v^2\frac{\dd}{\dd t}.
}
\tag{18}
\end{equation}

Hence
\begin{equation}
u'=v^2\dot u,
\qquad
v'=v^2\dot v,
\end{equation}
where the prime denotes differentiation with respect to $\tau$.

From (15),
\begin{align}
v'
&=
v^2
\frac{a-bv-uv^2-Iv^3}{v}\\
&=
av-bv^2-uv^3-Iv^4.
\end{align}

From (17),
\begin{align}
u'
&=
v^2
\frac{
au-(d+bu)v-u(1+u)v^2+(c-Iu)v^3
}{v^2}\\
&=
au-(d+bu)v-u(1+u)v^2+(c-Iu)v^3.
\end{align}

Thus the compactified vector field in the $x$-chart is

\begin{equation}
\boxed{
\begin{aligned}
u'
&=
au-(d+bu)v-u(1+u)v^2+(c-Iu)v^3,
\\[1mm]
v'
&=
av-bv^2-uv^3-Iv^4.
\end{aligned}}
\tag{19}
\end{equation}

This system is regular at $v=0$.

% =================================================
\subsection{4.4 Dynamics on the Line at Infinity}
% =================================================

Setting
\begin{equation}
v=0
\end{equation}
in (19), we obtain
\begin{equation}
\boxed{
u'=au,
\qquad
v'=0.
}
\tag{20}
\end{equation}

Therefore the singularity condition is
\begin{equation}
au=0.
\end{equation}

For $a\neq0$,
\begin{equation}
\boxed{
u_*=0.
}
\tag{21}
\end{equation}

Thus the singular point at infinity in the $x$-chart is
\begin{equation}
\boxed{
p_x=(0,0).
}
\tag{22}
\end{equation}

Since
\begin{equation}
u=\frac{y}{x},
\end{equation}
the condition $u=0$ means
\begin{equation}
\frac{y}{x}\to0.
\end{equation}

Therefore
\begin{equation}
\boxed{
p_x\text{ represents the $x$-axis direction at infinity}.
}
\tag{23}
\end{equation}

In homogeneous projective coordinates,
\begin{equation}
\boxed{
p_x=[1:0:0].
}
\tag{24}
\end{equation}

% =================================================
\subsection{4.5 Jacobian at $p_x$}
% =================================================

Define
\begin{align}
f(u,v)
&=
au-(d+bu)v-u(1+u)v^2+(c-Iu)v^3,
\\
g(u,v)
&=
av-bv^2-uv^3-Iv^4.
\end{align}

The Jacobian matrix is
\begin{equation}
J_x(u,v)
=
\begin{pmatrix}
f_u & f_v\\
g_u & g_v
\end{pmatrix}.
\end{equation}

We have
\begin{align}
f_u
&=
a-bv-(1+2u)v^2-Iv^3,
\\
f_v
&=
-(d+bu)-2u(1+u)v+3(c-Iu)v^2,
\\
g_u
&=
-v^3,
\\
g_v
&=
a-2bv-3uv^2-4Iv^3.
\end{align}

Therefore
\begin{equation}
\boxed{
J_x(u,v)
=
\begin{pmatrix}
a-bv-(1+2u)v^2-Iv^3
&
-(d+bu)-2u(1+u)v+3(c-Iu)v^2
\\[2mm]
-v^3
&
a-2bv-3uv^2-4Iv^3
\end{pmatrix}.
}
\tag{25}
\end{equation}

At $p_x=(0,0)$,
\begin{equation}
\boxed{
J_x(p_x)
=
\begin{pmatrix}
a & -d\\
0 & a
\end{pmatrix}.
}
\tag{26}
\end{equation}

The characteristic equation is
\begin{align}
\det
\begin{pmatrix}
a-\lambda & -d\\
0 & a-\lambda
\end{pmatrix}
&=0,
\end{align}
which gives
\begin{equation}
(a-\lambda)^2=0.
\end{equation}

Hence
\begin{equation}
\boxed{
\lambda_1=\lambda_2=a.
}
\tag{27}
\end{equation}

% =================================================
\subsection{4.6 Classification of $p_x$}
% =================================================

If
\begin{equation}
a>0,
\end{equation}
then
\begin{equation}
\lambda_1=\lambda_2=a>0,
\end{equation}
so $p_x$ is unstable.

For $d\neq0$, the matrix
\begin{equation}
\begin{pmatrix}
a&-d\\
0&a
\end{pmatrix}
\end{equation}
has only one eigenvector and therefore has a nontrivial Jordan block.
Consequently, $p_x$ is generically an unstable improper node.

Thus,
\begin{equation}
\boxed{
a>0,\ d\neq0
\quad\Longrightarrow\quad
p_x\text{ is an unstable improper node}.
}
\tag{28}
\end{equation}

If $d=0$, then
\begin{equation}
J_x(p_x)=aI,
\end{equation}
and the point is a star node.

Similarly, if
\begin{equation}
a<0,
\end{equation}
then
\begin{equation}
\lambda_1=\lambda_2=a<0,
\end{equation}
and $p_x$ is stable.

For $d\neq0$,
\begin{equation}
\boxed{
a<0,\ d\neq0
\quad\Longrightarrow\quad
p_x\text{ is a stable improper node}.
}
\tag{29}
\end{equation}

For $d=0$, it is a stable star node.

Therefore, for $a\neq0$, the point $p_x$ is hyperbolic.

% =================================================
\section{5. Second Local Chart: the $y$-Chart}
% =================================================

The first chart does not cover the vertical direction. We therefore
introduce the second local chart
\begin{equation}
\boxed{
u=\frac{x}{y},
\qquad
v=\frac1y.
}
\tag{30}
\end{equation}

The inverse transformation is
\begin{equation}
\boxed{
x=\frac{u}{v},
\qquad
y=\frac1v.
}
\tag{31}
\end{equation}

Again,
\begin{equation}
v=0
\end{equation}
represents infinity in this chart, while $u=x/y$ represents
direction.

% =================================================
\subsection{5.1 Transformation of the $v$-Equation}
% =================================================

Since
\begin{equation}
y=\frac1v,
\end{equation}
we have
\begin{equation}
\dot y=-\frac{\dot v}{v^2}.
\end{equation}

Since $\dot y=Q(x,y)$,
\begin{equation}
-\frac{\dot v}{v^2}
=
Q\left(\frac{u}{v},\frac1v\right),
\end{equation}
so
\begin{equation}
\boxed{
\dot v
=
-v^2Q\left(\frac{u}{v},\frac1v\right).
}
\tag{32}
\end{equation}

We have
\begin{align}
Q\left(\frac{u}{v},\frac1v\right)
&=
c-d\frac{u^2}{v^2}-\frac1v\\
&=
\frac{cv^3-du^2v-v^2}{v^3}.
\end{align}

Therefore
\begin{align}
\dot v
&=
-v^2
\frac{cv^3-du^2v-v^2}{v^3}\\
&=
\frac{du^2v+v^2-cv^3}{v}.
\end{align}

Hence
\begin{equation}
\boxed{
\dot v
=
\frac{du^2v+v^2-cv^3}{v}.
}
\tag{33}
\end{equation}

% =================================================
\subsection{5.2 Transformation of the $u$-Equation}
% =================================================

Since
\begin{equation}
u=\frac{x}{y},
\end{equation}
we have
\begin{equation}
\dot u
=
\frac{y\dot x-x\dot y}{y^2}.
\end{equation}

Using
\begin{equation}
x=\frac uv,\qquad
y=\frac1v,
\end{equation}
we obtain
\begin{equation}
\boxed{
\dot u
=
v\left[
P\left(\frac uv,\frac1v\right)
-uQ\left(\frac uv,\frac1v\right)
\right].
}
\tag{34}
\end{equation}

First,
\begin{align}
P\left(\frac uv,\frac1v\right)
&=
\frac1v
-a\frac{u^3}{v^3}
+b\frac{u^2}{v^2}
+I\\
&=
\frac{-au^3+bu^2v+uv^2+Iv^3}{v^3}.
\end{align}

Also,
\begin{align}
Q\left(\frac uv,\frac1v\right)
&=
\frac{cv^3-du^2v-v^2}{v^3}.
\end{align}

Therefore
\begin{align}
P-uQ
&=
\frac{
-au^3+(bu^2+du^3)v
+(u+u^2)v^2
+(I-cu)v^3
}{v^3}.
\end{align}

Consequently,
\begin{equation}
\boxed{
\dot u
=
\frac{
-au^3+(bu^2+du^3)v
+(u+u^2)v^2
+(I-cu)v^3
}{v^2}.
}
\tag{35}
\end{equation}

% =================================================
\subsection{5.3 Time Rescaling}
% =================================================

Again, since $n=3$,
\begin{equation}
\frac{\dd}{\dd\tau}
=
v^2\frac{\dd}{\dd t}.
\end{equation}

Thus
\begin{equation}
u'=v^2\dot u,
\qquad
v'=v^2\dot v.
\end{equation}

From (35),
\begin{align}
u'
&=
-au^3+(bu^2+du^3)v
+(u+u^2)v^2
+(I-cu)v^3.
\end{align}

From (33),
\begin{align}
v'
&=
v^2
\frac{du^2v+v^2-cv^3}{v}\\
&=
du^2v^2+v^3-cv^4.
\end{align}

Therefore the compactified system in the $y$-chart is

\begin{equation}
\boxed{
\begin{aligned}
u'
&=
-au^3+(bu^2+du^3)v
+(u+u^2)v^2
+(I-cu)v^3,
\\[1mm]
v'
&=
du^2v^2+v^3-cv^4.
\end{aligned}}
\tag{36}
\end{equation}

% =================================================
\subsection{5.4 Dynamics on the Line at Infinity}
% =================================================

Setting
\begin{equation}
v=0,
\end{equation}
we obtain
\begin{equation}
\boxed{
u'=-au^3,
\qquad
v'=0.
}
\tag{37}
\end{equation}

Hence the singularity condition is
\begin{equation}
-au^3=0.
\end{equation}

For $a\neq0$,
\begin{equation}
\boxed{
u_*=0.
}
\tag{38}
\end{equation}

Therefore the singular point in the $y$-chart is

\begin{equation}
\boxed{
p_y=(0,0).
}
\tag{39}
\end{equation}

Since
\begin{equation}
u=\frac{x}{y},
\end{equation}
the condition $u=0$ corresponds to
\begin{equation}
x=0.
\end{equation}

Hence

\begin{equation}
\boxed{
p_y\text{ represents the $y$-axis direction at infinity}.
}
\tag{40}
\end{equation}

In homogeneous projective coordinates,
\begin{equation}
\boxed{
p_y=[0:1:0].
}
\tag{41}
\end{equation}

% =================================================
\subsection{5.5 Jacobian at $p_y$}
% =================================================

Define
\begin{align}
f(u,v)
&=
-au^3+(bu^2+du^3)v
+(u+u^2)v^2
+(I-cu)v^3,
\\
g(u,v)
&=
du^2v^2+v^3-cv^4.
\end{align}

Every term in $f$ and $g$ has total degree at least $2$.
Consequently, there are no linear terms in the vector field at the
origin.

Hence
\begin{equation}
\boxed{
J_y(p_y)
=
\begin{pmatrix}
0&0\\
0&0
\end{pmatrix}.
}
\tag{42}
\end{equation}

Therefore the eigenvalues are
\begin{equation}
\boxed{
\lambda_1=\lambda_2=0.
}
\tag{43}
\end{equation}

Thus $p_y$ is a nonhyperbolic, degenerate singularity.

\begin{equation}
\boxed{
p_y\text{ is nonhyperbolic and cannot be classified by linearization alone}.
}
\tag{44}
\end{equation}

Higher-order terms must be considered in order to determine its local
phase portrait.

% =================================================
\section{6. Why the Two Infinity Points Are Different}
% =================================================

The singularity $p_x$ corresponds to
\begin{equation}
y=0,\qquad |x|\to\infty,
\end{equation}
whereas $p_y$ corresponds to
\begin{equation}
x=0,\qquad |y|\to\infty.
\end{equation}

Therefore they represent different directions at infinity:
\begin{equation}
\boxed{
p_x=[1:0:0],
\qquad
p_y=[0:1:0].
}
\end{equation}

The asymmetry between the two singularities is caused by the highest
degree part
\begin{equation}
(P_3,Q_3)=(-ax^3,0).
\end{equation}

At the $x$-axis direction,
\begin{equation}
y=0,
\end{equation}
and the cubic term
\begin{equation}
-ax^3
\end{equation}
is nonzero.

By contrast, along the $y$-axis direction,
\begin{equation}
x=0,
\end{equation}
and therefore
\begin{equation}
P_3(0,y)=0.
\end{equation}

This explains why the $x$-axis infinity point has a nonzero
linearization, whereas the $y$-axis infinity point is degenerate.

% =================================================
\section{7. Transition Between the Two Charts}
% =================================================

The two local charts are related by the transition map.

In the $x$-chart,
\begin{equation}
u_1=\frac{y}{x},
\qquad
v_1=\frac1x.
\end{equation}

In the $y$-chart,
\begin{equation}
u_2=\frac{x}{y},
\qquad
v_2=\frac1y.
\end{equation}

Therefore,
\begin{align}
u_2
&=
\frac{x}{y}
=
\frac{1}{y/x}
=
\frac1{u_1},
\\
v_2
&=
\frac1y
=
\frac{1/x}{y/x}
=
\frac{v_1}{u_1}.
\end{align}

Hence
\begin{equation}
\boxed{
u_2=\frac1{u_1},
\qquad
v_2=\frac{v_1}{u_1}.
}
\tag{45}
\end{equation}

This transition map is defined on the overlap where
\begin{equation}
u_1\neq0.
\end{equation}

The point $u_1=0$ is the horizontal direction, which is not
represented by the $y$-chart.

Similarly, $u_2=0$ represents the vertical direction, which is not
represented by the $x$-chart.

% =================================================
\section{8. Important Parameter Dependence}
% =================================================

The locations of the singular directions at infinity are determined
only by the highest-degree homogeneous terms
\begin{equation}
P_3(x,y)=-ax^3,
\qquad
Q_3(x,y)=0.
\end{equation}

Consequently, for $a\neq0$, the singular directions
\begin{equation}
[1:0:0],\qquad[0:1:0]
\end{equation}
do not depend on $b,c,d,$ or $I$.

However, the detailed local dynamics near the singularities may
depend on the lower-degree terms.

For example,
\begin{equation}
J_x(p_x)
=
\begin{pmatrix}
a&-d\\
0&a
\end{pmatrix},
\end{equation}
so $d$ affects the eigenvector and Jordan structure even though it
does not affect the location of the singular direction.

Thus the general principle is
\begin{equation}
\boxed{
\text{highest-degree terms}
\quad\Longrightarrow\quad
\text{singular directions at infinity},
}
\tag{46}
\end{equation}
whereas
\begin{equation}
\boxed{
\text{all polynomial terms}
\quad\Longrightarrow\quad
\text{local dynamics near infinity}.
}
\tag{47}
\end{equation}

% =================================================
\section{9. Degenerate Case $a=0$}
% =================================================

The preceding calculations assume
\begin{equation}
a\neq0.
\end{equation}

If
\begin{equation}
a=0,
\end{equation}
then the cubic term disappears and the system becomes
\begin{equation}
\dot x=y+bx^2+I,
\qquad
\dot y=c-dx^2-y.
\end{equation}

The degree is then generically
\begin{equation}
\boxed{
n=2,
}
\end{equation}
provided at least one quadratic coefficient is nonzero.

Consequently, the appropriate Poincaré time rescaling is no longer
$v^2$, but
\begin{equation}
\boxed{
\frac{\dd}{\dd\tau}
=
v\frac{\dd}{\dd t}.
}
\end{equation}

Thus, when parameters cause the degree of the polynomial vector field
to change, the compactification must be recomputed using the actual
degree of the resulting system.

% =================================================
\section{10. Reusable Procedure for a General Polynomial System}
% =================================================

For a general planar polynomial system
\begin{equation}
\dot x=P(x,y),
\qquad
\dot y=Q(x,y),
\end{equation}
the following procedure can be applied systematically.

\subsection*{Step 1: Determine the degree}

Set
\begin{equation}
n=\max\{\deg P,\deg Q\}.
\end{equation}

\subsection*{Step 2: Extract the highest-degree homogeneous terms}

Write
\begin{equation}
P=P_n+\cdots,
\qquad
Q=Q_n+\cdots.
\end{equation}

\subsection*{Step 3: Find the singular directions at infinity}

In the $x$-chart solve
\begin{equation}
\boxed{
Q_n(1,u)-uP_n(1,u)=0.
}
\end{equation}

In the $y$-chart solve
\begin{equation}
\boxed{
P_n(u,1)-uQ_n(u,1)=0.
}
\end{equation}

\subsection*{Step 4: Introduce the local charts}

For the $x$-chart,
\begin{equation}
\boxed{
x=\frac1v,
\qquad
y=\frac uv.
}
\end{equation}

For the $y$-chart,
\begin{equation}
\boxed{
x=\frac uv,
\qquad
y=\frac1v.
}
\end{equation}

\subsection*{Step 5: Transform the vector field}

In the $x$-chart,
\begin{equation}
\boxed{
\dot u
=
v(Q-uP),
\qquad
\dot v=-v^2P,
}
\end{equation}
where $P$ and $Q$ are evaluated at
\begin{equation}
\left(\frac1v,\frac uv\right).
\end{equation}

In the $y$-chart,
\begin{equation}
\boxed{
\dot u
=
v(P-uQ),
\qquad
\dot v=-v^2Q,
}
\end{equation}
where $P$ and $Q$ are evaluated at
\begin{equation}
\left(\frac uv,\frac1v\right).
\end{equation}

\subsection*{Step 6: Rescale time}

Use
\begin{equation}
\boxed{
\frac{\dd}{\dd\tau}
=
v^{n-1}\frac{\dd}{\dd t}.
}
\end{equation}

\subsection*{Step 7: Restrict to infinity}

Set
\begin{equation}
v=0.
\end{equation}

\subsection*{Step 8: Find and classify the singularities}

Solve the resulting equation for $u$, then calculate the Jacobian
matrix of the compactified vector field.

If the eigenvalues have nonzero real parts, the singularity is
hyperbolic and the local phase portrait is determined by the
linearization.

If zero eigenvalues occur, the singularity is nonhyperbolic and
higher-order analysis is required.

% =================================================
\section{11. Final Results for the Hindmarsh--Rose-Type System}
% =================================================

For the system
\begin{equation}
\boxed{
\begin{aligned}
\dot{x}&=y-a x^3+b x^2+I,\\
\dot{y}&=c-dx^2-y,
\end{aligned}}
\end{equation}
with $a\neq0$, the degree is
\begin{equation}
\boxed{n=3}.
\end{equation}

The highest-degree homogeneous part is
\begin{equation}
\boxed{
(P_3,Q_3)=(-ax^3,0).
}
\end{equation}

The singular directions at infinity are
\begin{equation}
\boxed{
[1:0:0]
\quad\text{and}\quad
[0:1:0].
}
\end{equation}

The point
\begin{equation}
\boxed{
p_x=[1:0:0]
}
\end{equation}
is represented in the $x$-chart by
\begin{equation}
(u,v)=(0,0)
\end{equation}
and has Jacobian
\begin{equation}
\boxed{
J_x(p_x)
=
\begin{pmatrix}
a&-d\\
0&a
\end{pmatrix}.
}
\end{equation}

Its eigenvalues are
\begin{equation}
\boxed{
\lambda_1=\lambda_2=a.
}
\end{equation}

Thus, for $a\neq0$,
\begin{equation}
\boxed{
\begin{cases}
a>0:& p_x\text{ is an unstable node},\\[1mm]
a<0:& p_x\text{ is a stable node}.
\end{cases}
}
\end{equation}

For $d\neq0$, the node is generically an improper node.

The point
\begin{equation}
\boxed{
p_y=[0:1:0]
}
\end{equation}
is represented in the $y$-chart by
\begin{equation}
(u,v)=(0,0)
\end{equation}
and has Jacobian
\begin{equation}
\boxed{
J_y(p_y)
=
\begin{pmatrix}
0&0\\
0&0
\end{pmatrix}.
}
\end{equation}

Hence
\begin{equation}
\boxed{
\lambda_1=\lambda_2=0,
}
\end{equation}
so
\begin{equation}
\boxed{
p_y\text{ is a degenerate nonhyperbolic singularity at infinity}.
}
\end{equation}

Its local phase portrait cannot be obtained from the linearization and
requires an analysis of the higher-order terms.

\section{1. Equilibrium Equations}


An equilibrium point
\begin{equation}
E^*=(x^*,y^*)
\end{equation}
satisfies
\begin{equation}
\dot{x}=0,
\qquad
\dot{y}=0.
\end{equation}

Therefore,
\begin{equation}
y^*-a(x^*)^3+b(x^*)^2+I=0,
\tag{1}
\end{equation}
and
\begin{equation}
c-d(x^*)^2-y^*=0.
\tag{2}
\end{equation}

From the second equation,
\begin{equation}
\boxed{
y^*=c-d(x^*)^2.
}
\tag{3}
\end{equation}

Substituting (3) into the first equilibrium equation gives
\begin{align}
c-d(x^*)^2-a(x^*)^3+b(x^*)^2+I
&=0.
\end{align}

Hence
\begin{equation}
-a(x^*)^3+(b-d)(x^*)^2+(c+I)=0.
\end{equation}

Multiplying by $-1$ gives
\begin{equation}
\boxed{
a(x^*)^3+(d-b)(x^*)^2-(c+I)=0.
}
\tag{4}
\end{equation}

Thus all finite equilibria are obtained by solving the cubic
equation (4), and then using (3) to obtain their $y$-coordinates.

For convenience, define
\begin{equation}
\boxed{
B=d-b,
\qquad
S=c+I.
}
\tag{5}
\end{equation}

Then the equilibrium equation becomes
\begin{equation}
\boxed{
a x^3+B x^2-S=0.
}
\tag{6}
\end{equation}

Hence every equilibrium is of the form
\begin{equation}
\boxed{
E^*=
\left(
x^*,
c-d(x^*)^2
\right),
}
\tag{7}
\end{equation}
where $x^*$ is a real root of (6).

% =========================================================
\section{2. Number of Equilibria}
% =========================================================

Assume
\begin{equation}
a\neq0.
\end{equation}

Then (6) is a cubic equation. Its discriminant is
\begin{equation}
\boxed{
\mathcal{D}
=
4B^3S-27a^2S^2.
}
\tag{8}
\end{equation}

Returning to the original parameters,
\begin{equation}
\boxed{
\mathcal{D}
=
4(d-b)^3(c+I)
-
27a^2(c+I)^2.
}
\tag{9}
\end{equation}

The number of distinct real equilibria is determined by the sign of
$\mathcal{D}$:

\begin{equation}
\boxed{
\begin{array}{c|c}
\mathcal{D}>0
&
\text{three distinct real equilibria},
\\[2mm]
\mathcal{D}=0
&
\text{at least one multiple real equilibrium},
\\[2mm]
\mathcal{D}<0
&
\text{one real equilibrium}.
\end{array}}
\tag{10}
\end{equation}

Thus the discriminant provides a direct criterion for the number of
finite equilibria.

% =========================================================
\section{3. Explicit Solution of the Cubic}
% =========================================================

Starting from
\begin{equation}
a x^3+B x^2-S=0,
\end{equation}
divide by $a$:
\begin{equation}
x^3+\frac{B}{a}x^2-\frac{S}{a}=0.
\tag{11}
\end{equation}

Introduce the change of variable
\begin{equation}
\boxed{
x=z-\frac{B}{3a}.
}
\tag{12}
\end{equation}

This removes the quadratic term and yields the depressed cubic
\begin{equation}
z^3+pz+q=0,
\tag{13}
\end{equation}
where
\begin{equation}
\boxed{
p=-\frac{B^2}{3a^2},
}
\tag{14}
\end{equation}
and
\begin{equation}
\boxed{
q=
\frac{2B^3}{27a^3}
-\frac{S}{a}.
}
\tag{15}
\end{equation}

The Cardano quantity is
\begin{equation}
\Delta
=
\left(\frac q2\right)^2
+
\left(\frac p3\right)^3.
\tag{16}
\end{equation}

After simplification,
\begin{equation}
\boxed{
\Delta
=
\frac{
27a^2S^2-4B^3S
}{
108a^4
}.
}
\tag{17}
\end{equation}

In terms of the cubic discriminant,
\begin{equation}
\boxed{
\mathcal{D}
=
-108a^4\Delta.
}
\tag{18}
\end{equation}

Therefore,
\begin{equation}
\mathcal{D}>0
\quad\Longleftrightarrow\quad
\Delta<0,
\end{equation}
which corresponds to three distinct real roots.

% =========================================================
\section{4. The One-Real-Equilibrium Case}
% =========================================================

When
\begin{equation}
\Delta>0
\end{equation}
or equivalently
\begin{equation}
\mathcal{D}<0,
\end{equation}
there is one real root.

By Cardano's formula,
\begin{equation}
z=
\sqrt[3]{-\frac q2+\sqrt{\Delta}}
+
\sqrt[3]{-\frac q2-\sqrt{\Delta}}.
\end{equation}

Consequently,
\begin{equation}
\boxed{
x^*
=
-\frac{B}{3a}
+
\sqrt[3]{-\frac q2+\sqrt{\Delta}}
+
\sqrt[3]{-\frac q2-\sqrt{\Delta}}.
}
\tag{19}
\end{equation}

Using
\begin{equation}
-\frac q2
=
\frac{S}{2a}
-\frac{B^3}{27a^3},
\end{equation}
we obtain
\begin{equation}
\boxed{
x^*
=
-\frac{B}{3a}
+
\sqrt[3]{
\frac{S}{2a}
-\frac{B^3}{27a^3}
+\sqrt{\Delta}}
+
\sqrt[3]{
\frac{S}{2a}
-\frac{B^3}{27a^3}
-\sqrt{\Delta}}
.
}
\tag{20}
\end{equation}

The corresponding equilibrium is
\begin{equation}
\boxed{
E^*
=
\left(
x^*,
c-d(x^*)^2
\right).
}
\tag{21}
\end{equation}

% =========================================================
\section{5. The Three-Real-Equilibria Case}
% =========================================================

When
\begin{equation}
\mathcal{D}>0,
\end{equation}
the cubic has three distinct real roots.

Since
\begin{equation}
p=-\frac{B^2}{3a^2}<0
\end{equation}
for $B\neq0$, the trigonometric form of the cubic solution is
convenient.

Define
\begin{equation}
\boxed{
\theta
=
\frac13
\arccos
\left(
1-\frac{27a^2S}{2B^3}
\right).
}
\tag{22}
\end{equation}

Then the three real roots may be written as
\begin{equation}
\boxed{
x_k
=
-\frac{B}{3a}
+
\frac{2|B|}{3|a|}
\cos\left(
\theta-\frac{2\pi k}{3}
\right),
\qquad
k=0,1,2.
}
\tag{23}
\end{equation}

The corresponding equilibria are
\begin{equation}
\boxed{
E_k=
\left(
x_k,
c-dx_k^2
\right),
\qquad
k=0,1,2.
}
\tag{24}
\end{equation}

For numerical calculations, it is often more convenient to obtain
the roots using the numerical polynomial-root algorithm rather than
the explicit trigonometric formula. Nevertheless, formula (23) gives
an analytical representation of all three roots.

% =========================================================
\section{6. Special Case $c+I=0$}
% =========================================================

Suppose
\begin{equation}
S=c+I=0.
\end{equation}

Then the equilibrium equation becomes
\begin{equation}
ax^3+Bx^2=0.
\end{equation}

Factoring,
\begin{equation}
x^2(ax+B)=0.
\end{equation}

Hence
\begin{equation}
\boxed{
x_1^*=0
}
\tag{25}
\end{equation}
is a double root, while
\begin{equation}
\boxed{
x_2^*=-\frac{B}{a}
=
\frac{b-d}{a}.
}
\tag{26}
\end{equation}

The corresponding $y$-coordinates are
\begin{equation}
y_1^*=c,
\end{equation}
and
\begin{equation}
y_2^*
=
c-d
\left(
\frac{b-d}{a}
\right)^2.
\end{equation}

Therefore the distinct equilibria are
\begin{equation}
\boxed{
E_1=(0,c),
}
\tag{27}
\end{equation}
and
\begin{equation}
\boxed{
E_2=
\left(
\frac{b-d}{a},
\,
c-d\frac{(b-d)^2}{a^2}
\right).
}
\tag{28}
\end{equation}

The root $x=0$ is algebraically double and therefore corresponds to a
degenerate parameter configuration.

% =========================================================
\section{7. Jacobian Matrix at an Arbitrary Equilibrium}
% =========================================================

The vector field is
\begin{equation}
F(x,y)
=
\begin{pmatrix}
y-a x^3+b x^2+I\\
c-dx^2-y
\end{pmatrix}.
\end{equation}

The Jacobian matrix is
\begin{equation}
DF(x,y)
=
\begin{pmatrix}
\dfrac{\partial\dot x}{\partial x}
&
\dfrac{\partial\dot x}{\partial y}
\\[3mm]
\dfrac{\partial\dot y}{\partial x}
&
\dfrac{\partial\dot y}{\partial y}
\end{pmatrix}.
\end{equation}

Hence
\begin{equation}
\boxed{
DF(x,y)
=
\begin{pmatrix}
-3ax^2+2bx&1\\
-2dx&-1
\end{pmatrix}.
}
\tag{29}
\end{equation}

At an equilibrium
\begin{equation}
E^*=(x^*,y^*),
\end{equation}
the Jacobian is therefore
\begin{equation}
\boxed{
J^*
=
\begin{pmatrix}
-3a(x^*)^2+2bx^*&1\\
-2dx^*&-1
\end{pmatrix}.
}
\tag{30}
\end{equation}

% =========================================================
\section{8. Trace and Determinant}
% =========================================================

The trace of $J^*$ is
\begin{align}
\tr(J^*)
&=
-3a(x^*)^2+2bx^*-1.
\end{align}

Thus
\begin{equation}
\boxed{
T(x^*)
=
-3a(x^*)^2+2bx^*-1.
}
\tag{31}
\end{equation}

The determinant is
\begin{align}
\detm(J^*)
&=
\left[-3a(x^*)^2+2bx^*\right](-1)
-
(1)(-2dx^*)\\
&=
3a(x^*)^2-2bx^*+2dx^*.
\end{align}

Therefore
\begin{equation}
\boxed{
D(x^*)
=
3a(x^*)^2+2(d-b)x^*.
}
\tag{32}
\end{equation}

Using $B=d-b$,
\begin{equation}
\boxed{
D(x^*)
=
x^*\left(3ax^*+2B\right).
}
\tag{33}
\end{equation}

% =========================================================
\section{9. Eigenvalues}
% =========================================================

The characteristic polynomial of $J^*$ is
\begin{equation}
\lambda^2-T\lambda+D=0.
\tag{34}
\end{equation}

Consequently,
\begin{equation}
\boxed{
\lambda_{\pm}
=
\frac12
\left[
T\pm\sqrt{T^2-4D}
\right].
}
\tag{35}
\end{equation}

Here
\begin{equation}
T=-3a(x^*)^2+2bx^*-1
\end{equation}
and
\begin{equation}
D=
x^*
\left[
3ax^*+2(d-b)
\right].
\end{equation}

% =========================================================
\section{10. Classification of Equilibria}
% =========================================================

The local type of an equilibrium can be determined from its trace,
determinant, and discriminant.

\subsection*{10.1 Saddle}

If
\begin{equation}
\boxed{
D<0,
}
\tag{36}
\end{equation}
then the eigenvalues have opposite signs. Therefore
\begin{equation}
\boxed{
E^*\text{ is a saddle}.
}
\tag{37}
\end{equation}

\subsection*{10.2 Stable Equilibrium}

If
\begin{equation}
D>0
\end{equation}
and
\begin{equation}
T<0,
\end{equation}
then both eigenvalues have negative real parts.

Therefore
\begin{equation}
\boxed{
E^*\text{ is locally asymptotically stable}.
}
\tag{38}
\end{equation}

More precisely:

\begin{equation}
T^2-4D>0
\quad\Longrightarrow\quad
\text{stable node},
\end{equation}

whereas

\begin{equation}
T^2-4D<0
\quad\Longrightarrow\quad
\text{stable focus}.
\end{equation}

\subsection*{10.3 Unstable Equilibrium}

If
\begin{equation}
D>0
\end{equation}
and
\begin{equation}
T>0,
\end{equation}
then both eigenvalues have positive real parts.

Hence
\begin{equation}
\boxed{
E^*\text{ is unstable}.
}
\tag{39}
\end{equation}

More precisely,

\begin{equation}
T^2-4D>0
\quad\Longrightarrow\quad
\text{unstable node},
\end{equation}

while

\begin{equation}
T^2-4D<0
\quad\Longrightarrow\quad
\text{unstable focus}.
\end{equation}

% =========================================================
\section{11. Hopf Bifurcation Condition}
% =========================================================

For a planar system, a Hopf bifurcation requires a pair of purely
imaginary eigenvalues. For the characteristic polynomial
\begin{equation}
\lambda^2-T\lambda+D=0,
\end{equation}
the necessary linear conditions are
\begin{equation}
\boxed{
T=0,
\qquad
D>0.
}
\tag{40}
\end{equation}

For the present system,
\begin{equation}
T(x^*)
=
-3a(x^*)^2+2bx^*-1.
\end{equation}

Thus the trace condition is
\begin{equation}
-3a(x^*)^2+2bx^*-1=0,
\end{equation}
or
\begin{equation}
\boxed{
3a(x^*)^2-2bx^*+1=0.
}
\tag{41}
\end{equation}

This equation determines the candidate equilibrium locations at Hopf.

% =========================================================
\section{12. Hopf Equilibrium Coordinates}
% =========================================================

Solving
\begin{equation}
3a(x^*)^2-2bx^*+1=0
\end{equation}
gives
\begin{equation}
\boxed{
x_H^{\pm}
=
\frac{
b\pm\sqrt{b^2-3a}
}{
3a
}.
}
\tag{42}
\end{equation}

Therefore real Hopf candidates require
\begin{equation}
\boxed{
b^2-3a\ge0.
}
\tag{43}
\end{equation}

For each candidate $x_H^\pm$, the corresponding equilibrium
$y$-coordinate is
\begin{equation}
\boxed{
y_H^\pm
=
c-d(x_H^\pm)^2.
}
\tag{44}
\end{equation}

Thus the candidate Hopf equilibria are
\begin{equation}
\boxed{
E_H^\pm
=
\left(
x_H^\pm,
c-d(x_H^\pm)^2
\right).
}
\tag{45}
\end{equation}

% =========================================================
\section{13. Determinant at the Hopf Point}
% =========================================================

At a Hopf candidate,
\begin{equation}
x_H^\pm
=
\frac{
b\pm\sqrt{b^2-3a}
}{
3a
}.
\end{equation}

The determinant is
\begin{equation}
D_H
=
x_H
\left[
3ax_H+2(d-b)
\right].
\tag{46}
\end{equation}

Since
\begin{equation}
3ax_H^\pm
=
b\pm\sqrt{b^2-3a},
\end{equation}
we obtain
\begin{align}
3ax_H^\pm+2(d-b)
&=
b\pm\sqrt{b^2-3a}+2d-2b\\
&=
2d-b\pm\sqrt{b^2-3a}.
\end{align}

Therefore
\begin{equation}
\boxed{
D_H^\pm
=
x_H^\pm
\left[
2d-b\pm\sqrt{b^2-3a}
\right].
}
\tag{47}
\end{equation}

For a Hopf bifurcation candidate, we require
\begin{equation}
\boxed{
D_H^\pm>0.
}
\tag{48}
\end{equation}

At such a point,
\begin{equation}
\lambda_{\pm}
=
\pm i\omega_H,
\end{equation}
where
\begin{equation}
\boxed{
\omega_H=\sqrt{D_H}.
}
\tag{49}
\end{equation}

% =========================================================
\section{14. Hopf Parameter Condition}
% =========================================================

The equilibrium equation is
\begin{equation}
a x^3+(d-b)x^2-(c+I)=0.
\end{equation}

Therefore, at a Hopf point $x=x_H$, we must have
\begin{equation}
a x_H^3+(d-b)x_H^2-(c+I_H)=0.
\end{equation}

Hence
\begin{equation}
\boxed{
c+I_H
=
a x_H^3+(d-b)x_H^2.
}
\tag{50}
\end{equation}

Consequently,
\begin{equation}
\boxed{
I_H
=
a x_H^3+(d-b)x_H^2-c.
}
\tag{51}
\end{equation}
 
For the two Hopf candidates $x_H^\pm$,
\begin{equation}
\boxed{
I_H^\pm
=
a(x_H^\pm)^3
+
(d-b)(x_H^\pm)^2
-c.
}
\tag{52}
\end{equation}

Thus the Hopf bifurcation candidates can be parameterized as
\begin{equation}
\boxed{
x_H^\pm
=
\frac{
b\pm\sqrt{b^2-3a}
}{
3a
},
}
\tag{53}
\end{equation}
and
\begin{equation}
\boxed{
I_H^\pm
=
a(x_H^\pm)^3
+
(d-b)(x_H^\pm)^2
-c.
}
\tag{54}
\end{equation}

% =========================================================
\section{15. Simplification of the Hopf Parameter Relation}
% =========================================================

At a Hopf point,
\begin{equation}
3ax_H^2-2bx_H+1=0.
\end{equation}

Hence
\begin{equation}
3ax_H^2=2bx_H-1.
\tag{55}
\end{equation}

Multiplying by $x_H/3$,
\begin{equation}
a x_H^3
=
\frac{x_H(2bx_H-1)}{3}.
\tag{56}
\end{equation}

Therefore
\begin{align}
c+I_H
&=
a x_H^3+(d-b)x_H^2\\
&=
\frac{2b}{3}x_H^2-\frac{x_H}{3}
+(d-b)x_H^2\\
&=
\left(d-\frac b3\right)x_H^2
-\frac{x_H}{3}.
\end{align}

Consequently,
\begin{equation}
\boxed{
I_H
=
\left(d-\frac b3\right)x_H^2
-\frac{x_H}{3}
-c.
}
\tag{57}
\end{equation}

This form is often convenient for algebraic simplification.

% =========================================================
\section{16. Relation to Limit Cycles}
% =========================================================

The existence of a Hopf candidate does not by itself prove the
existence of a periodic orbit.

The logical sequence is
\begin{equation}
\boxed{
\text{equilibrium}
\quad\longrightarrow\quad
T=0,\ D>0
\quad\longrightarrow\quad
\text{Hopf candidate}.
}
\tag{58}
\end{equation}

For a generic Hopf bifurcation, one additionally needs the appropriate
nondegeneracy and transversality conditions.

In particular, the eigenvalues must cross the imaginary axis as the
bifurcation parameter varies, and the first Lyapunov coefficient must
satisfy
\begin{equation}
\boxed{
l_1\neq0.
}
\tag{59}
\end{equation}

When the generic Hopf hypotheses are satisfied, a small-amplitude
limit cycle is created or destroyed at the Hopf bifurcation.

Thus
\begin{equation}
\boxed{
T=0,\ D>0
}
\end{equation}
provides only the linear Hopf condition, whereas
\begin{equation}
\boxed{
l_1\neq0
}
\end{equation}
is needed to establish a generic nondegenerate Hopf bifurcation.

% =========================================================
\section{17. Divergence and a No-Limit-Cycle Criterion}
% =========================================================

The divergence of the vector field is
\begin{align}
\nabla\cdot F
&=
\frac{\partial \dot{x}}{\partial x}
+
\frac{\partial \dot{y}}{\partial y}\\
&=
(-3ax^2+2bx)+(-1).
\end{align}

Therefore
\begin{equation}
\boxed{
\nabla\cdot F
=
-3ax^2+2bx-1.
}
\tag{60}
\end{equation}

Notice that this is precisely the trace of the Jacobian:
\begin{equation}
\boxed{
\nabla\cdot F=T(x).
}
\tag{61}
\end{equation}

Assume
\begin{equation}
a>0.
\end{equation}

The divergence is then a concave quadratic function of $x$:
\begin{equation}
\nabla\cdot F
=
-3ax^2+2bx-1.
\end{equation}

Its maximum occurs at
\begin{equation}
x=\frac{b}{3a}.
\end{equation}

The maximum value is
\begin{align}
\max_x\nabla\cdot F
&=
-3a\left(\frac{b}{3a}\right)^2
+
2b\left(\frac{b}{3a}\right)-1\\
&=
-\frac{b^2}{3a}
+
\frac{2b^2}{3a}
-1\\
&=
\frac{b^2}{3a}-1\\
&=
\frac{b^2-3a}{3a}.
\end{align}

Consequently, if
\begin{equation}
\boxed{
a>0,
\qquad
b^2\le3a,
}
\tag{62}
\end{equation}
then
\begin{equation}
\boxed{
\nabla\cdot F\le0
\qquad\text{for all }(x,y).
}
\tag{63}
\end{equation}

Since the divergence is not identically zero, Bendixson's criterion
rules out nonconstant periodic orbits in the plane.

Therefore,
\begin{equation}
\boxed{
a>0,\qquad b^2\le3a
\quad\Longrightarrow\quad
\text{no planar limit cycle}.
}
\tag{64}
\end{equation}

This result is also consistent with the Hopf condition. Indeed, for
\begin{equation}
a>0,\qquad b^2<3a,
\end{equation}
the equation
\begin{equation}
3ax^2-2bx+1=0
\end{equation}
has no real solution, so no Hopf candidate exists.

% =========================================================
\section{18. Practical Procedure for a Concrete Parameter Set}
% =========================================================

Suppose that numerical values of
\begin{equation}
(a,b,c,d,I)
\end{equation}
are specified.

The procedure for investigating equilibria and possible limit cycles
is as follows.

\subsection*{Step 1: Solve the equilibrium cubic}

Solve
\begin{equation}
\boxed{
a x^3+(d-b)x^2-(c+I)=0.
}
\tag{65}
\end{equation}

Let the real roots be
\begin{equation}
x_1^*,x_2^*,x_3^*,
\end{equation}
where only the real roots are retained.

\subsection*{Step 2: Calculate the corresponding $y$-coordinates}

For each real root,
\begin{equation}
\boxed{
y_i^*
=
c-d(x_i^*)^2.
}
\tag{66}
\end{equation}

Hence
\begin{equation}
\boxed{
E_i=
\left(
x_i^*,
c-d(x_i^*)^2
\right).
}
\tag{67}
\end{equation}

\subsection*{Step 3: Calculate the trace}

For each equilibrium,
\begin{equation}
\boxed{
T_i
=
-3a(x_i^*)^2+2bx_i^*-1.
}
\tag{68}
\end{equation}

\subsection*{Step 4: Calculate the determinant}

For each equilibrium,
\begin{equation}
\boxed{
D_i
=
x_i^*
\left[
3ax_i^*+2(d-b)
\right].
}
\tag{69}
\end{equation}

\subsection*{Step 5: Classify the equilibrium}

If
\begin{equation}
D_i<0,
\end{equation}
then $E_i$ is a saddle.

If
\begin{equation}
D_i>0,\qquad T_i<0,
\end{equation}
then $E_i$ is stable.

If
\begin{equation}
D_i>0,\qquad T_i>0,
\end{equation}
then $E_i$ is unstable.

\subsection*{Step 6: Search for Hopf points}

Solve
\begin{equation}
\boxed{
3ax^2-2bx+1=0.
}
\tag{70}
\end{equation}

The candidate values are
\begin{equation}
\boxed{
x_H^\pm
=
\frac{b\pm\sqrt{b^2-3a}}{3a}.
}
\tag{71}
\end{equation}

Retain only those candidates for which
\begin{equation}
\boxed{
D_H^\pm>0.
}
\tag{72}
\end{equation}

Then calculate
\begin{equation}
\boxed{
I_H^\pm
=
a(x_H^\pm)^3
+
(d-b)(x_H^\pm)^2
-c.
}
\tag{73}
\end{equation}

Finally, to establish a genuine generic Hopf bifurcation, calculate
the first Lyapunov coefficient and verify
\begin{equation}
\boxed{
l_1\neq0.
}
\tag{74}
\end{equation}

% =========================================================
\section{19. MATLAB Implementation}
% =========================================================

The following MATLAB code computes the finite equilibria, their
Jacobians, eigenvalues, trace, determinant, and basic classification.

\begin{verbatim}
clear;
clc;

%% Parameters
a = 1;
b = 2;
c = 1;
d = 1;
I = 1;

%% Equilibrium equation
% a*x^3 + (d-b)*x^2 - (c+I) = 0

coef = [a, d-b, 0, -(c+I)];

xroots = roots(coef);

tol = 1e-10;

fprintf('============================================\n');
fprintf('Finite equilibria\n');
fprintf('============================================\n');

count = 0;

for k = 1:length(xroots)

    if abs(imag(xroots(k))) < tol

        count = count + 1;

        xstar = real(xroots(k));
        ystar = c - d*xstar^2;

        %% Jacobian
        J = [ -3*a*xstar^2 + 2*b*xstar,  1;
              -2*d*xstar,               -1 ];

        %% Eigenvalues
        eigJ = eig(J);

        %% Trace and determinant
        T = trace(J);
        D = det(J);

        fprintf('\nEquilibrium %d\n', count);
        fprintf('x* = %.12f\n', xstar);
        fprintf('y* = %.12f\n', ystar);

        fprintf('Trace = %.12f\n', T);
        fprintf('Determinant = %.12f\n', D);

        fprintf('Eigenvalues:\n');
        disp(eigJ);

        %% Classification
        if D < 0

            fprintf('Classification: Saddle\n');

        elseif D > 0 && T < 0

            if T^2 - 4*D > 0
                fprintf('Classification: Stable node\n');
            elseif T^2 - 4*D < 0
                fprintf('Classification: Stable focus\n');
            else
                fprintf('Classification: Degenerate stable equilibrium\n');
            end

        elseif D > 0 && T > 0

            if T^2 - 4*D > 0
                fprintf('Classification: Unstable node\n');
            elseif T^2 - 4*D < 0
                fprintf('Classification: Unstable focus\n');
            else
                fprintf('Classification: Degenerate unstable equilibrium\n');
            end

        elseif D > 0 && abs(T) < tol

            fprintf('Hopf candidate: T = 0 and D > 0\n');
            fprintf('Hopf frequency = %.12f\n', sqrt(D));

        else

            fprintf('Classification: Nonhyperbolic / degenerate\n');

        end
    end
end

%% Cubic discriminant
B = d - b;
S = c + I;

Disc = 4*B^3*S - 27*a^2*S^2;

fprintf('\n============================================\n');
fprintf('Cubic discriminant\n');
fprintf('============================================\n');

fprintf('Discriminant = %.12f\n', Disc);

if Disc > tol
    fprintf('Three distinct real equilibria.\n');
elseif abs(Disc) <= tol
    fprintf('Multiple equilibrium(s).\n');
else
    fprintf('One real equilibrium.\n');
end

%% Hopf candidates
fprintf('\n============================================\n');
fprintf('Hopf candidates\n');
fprintf('============================================\n');

if b^2 - 3*a >= 0

    xH1 = (b + sqrt(b^2 - 3*a))/(3*a);
    xH2 = (b - sqrt(b^2 - 3*a))/(3*a);

    xH = [xH1, xH2];

    for k = 1:2

        xh = xH(k);

        DH = xh*(3*a*xh + 2*(d-b));

        IH = a*xh^3 + (d-b)*xh^2 - c;

        fprintf('\nHopf candidate %d\n', k);
        fprintf('x_H = %.12f\n', xh);
        fprintf('D_H = %.12f\n', DH);
        fprintf('I_H = %.12f\n', IH);

        if DH > 0
            fprintf('Linear Hopf condition satisfied: D_H > 0.\n');
            fprintf('omega_H = %.12f\n', sqrt(DH));
        else
            fprintf('Not a Hopf point because D_H <= 0.\n');
        end
    end

else

    fprintf('No real Hopf candidates because b^2 - 3*a < 0.\n');

end

%% Divergence test
if a > 0

    maxDiv = (b^2 - 3*a)/(3*a);

    fprintf('\n============================================\n');
    fprintf('Divergence criterion\n');
    fprintf('============================================\n');

    fprintf('Maximum divergence = %.12f\n', maxDiv);

    if maxDiv <= tol
        fprintf('Divergence <= 0 for all x.\n');
        fprintf('No planar limit cycle by Bendixson criterion.\n');
    else
        fprintf('Divergence changes sign or can be positive.\n');
        fprintf('Bendixson criterion does not rule out limit cycles.\n');
    end

end
\end{verbatim}

% =========================================================
\section*{Summary of the Main Formulas}
% =========================================================

For the two-dimensional Hindmarsh--Rose-type model
\begin{equation}
\boxed{
\begin{aligned}
\dot{x}&=y-a x^3+b x^2+I,\\
\dot{y}&=c-dx^2-y,
\end{aligned}}
\end{equation}
the main formulas for the finite-equilibrium and limit-cycle analysis
are summarized below.

\subsection*{Equilibrium equation}

\begin{equation}
\boxed{
a x^3+(d-b)x^2-(c+I)=0.
}
\end{equation}

\subsection*{$y$-coordinate}

\begin{equation}
\boxed{
y^*=c-d(x^*)^2.
}
\end{equation}

\subsection*{Cubic discriminant}

\begin{equation}
\boxed{
\mathcal{D}
=
4(d-b)^3(c+I)
-
27a^2(c+I)^2.
}
\end{equation}

\subsection*{Jacobian}

\begin{equation}
\boxed{
J^*
=
\begin{pmatrix}
-3a(x^*)^2+2bx^*&1\\
-2dx^*&-1
\end{pmatrix}.
}
\end{equation}

\subsection*{Trace}

\begin{equation}
\boxed{
T
=
-3a(x^*)^2+2bx^*-1.
}
\end{equation}

\subsection*{Determinant}

\begin{equation}
\boxed{
D
=
x^*
\left[
3ax^*+2(d-b)
\right].
}
\end{equation}

\subsection*{Eigenvalues}

\begin{equation}
\boxed{
\lambda_{\pm}
=
\frac12
\left[
T\pm\sqrt{T^2-4D}
\right].
}
\end{equation}

\subsection*{Hopf condition}

\begin{equation}
\boxed{
T=0,\qquad D>0.
}
\end{equation}

Equivalently,

\begin{equation}
\boxed{
3ax_H^2-2bx_H+1=0.
}
\end{equation}

\subsection*{Hopf equilibrium coordinates}

\begin{equation}
\boxed{
x_H^\pm
=
\frac{
b\pm\sqrt{b^2-3a}
}{
3a
}.
}
\end{equation}

\subsection*{Hopf parameter values}

\begin{equation}
\boxed{
I_H^\pm
=
a(x_H^\pm)^3
+
(d-b)(x_H^\pm)^2
-c.
}
\end{equation}

\subsection*{Hopf frequency}

\begin{equation}
\boxed{
\omega_H=\sqrt{D_H}.
}
\end{equation}

\subsection*{No-limit-cycle condition}

For
\begin{equation}
a>0,\qquad b^2\le3a,
\end{equation}
we have
\begin{equation}
\boxed{
\nabla\cdot F\le0
\quad\text{everywhere},
}
\end{equation}
and therefore, by Bendixson's criterion,

\begin{equation}
\boxed{
\text{there is no nonconstant periodic orbit in the plane}.
}
\end{equation}


oooooooooooooooo

% ============================================================
\section{The two-dimensional Hindmarsh--Rose model and its
reduced parameterization}
% ============================================================

We begin with the two-dimensional Hindmarsh--Rose-type system
\begin{equation}
\boxed{
\begin{aligned}
\dot{x} &= y-a x^3+b x^2+I,\\
\dot{y} &= c-dx^2-y,
\end{aligned}}
\tag{1}
\end{equation}
where $x$ and $y$ denote the fast state variables and
$a,b,c,d,I$ are real parameters. In the neuronal interpretation,
$x$ represents the membrane potential, while $y$ is a fast recovery
(or gating) variable. The parameter $I$ represents the externally
applied input current, whereas $a,b,c,$ and $d$ are intrinsic
parameters governing the nonlinear activation and recovery
mechanisms.

The system in \eqref{1} contains five parameters. Although retaining
all five parameters is useful when describing the general model, a
direct analytical treatment of the full five-parameter family is
unnecessarily cumbersome for the qualitative analysis pursued in
this work, particularly for the determination of equilibria, Hopf
bifurcations, periodic orbits, and the global phase portrait.
Accordingly, we introduce a mathematically justified reduction of the
parameter space while preserving the main biological interpretation
of the model.

% ------------------------------------------------------------
\subsection{Normalization of the cubic coefficient}
% ------------------------------------------------------------

We first note that, for the case
\begin{equation}
a>0,
\end{equation}
the coefficient $a$ can be normalized by an appropriate rescaling of
the state variables. Indeed, introduce
\begin{equation}
X=\sqrt{a}\,x,
\qquad
Y=\sqrt{a}\,(y-c).
\tag{2}
\end{equation}
Equivalently,
\begin{equation}
x=\frac{X}{\sqrt{a}},
\qquad
y=c+\frac{Y}{\sqrt{a}}.
\end{equation}
Since the transformation is a smooth invertible change of
coordinates for $a>0$, it preserves the qualitative phase-space
structure up to a smooth coordinate transformation.

Substituting \eqref{2} into \eqref{1}, we obtain
\begin{equation}
\dot{X}
=
Y-X^3+\frac{b}{\sqrt{a}}X^2+\sqrt{a}(c+I),
\tag{3}
\end{equation}
and
\begin{equation}
\dot{Y}
=
-\frac{d}{\sqrt{a}}X^2-Y.
\tag{4}
\end{equation}

Thus, if we introduce the effective parameters
\begin{equation}
\beta=\frac{b}{\sqrt{a}},
\qquad
\delta=\frac{d}{\sqrt{a}},
\qquad
\mu=\sqrt{a}(c+I),
\tag{5}
\end{equation}
the system takes the normalized form
\begin{equation}
\boxed{
\begin{aligned}
\dot{X}&=Y-X^3+\beta X^2+\mu,\\
\dot{Y}&=-\delta X^2-Y.
\end{aligned}}
\tag{6}
\end{equation}

This transformation shows that, for the qualitative analysis of the
two-dimensional autonomous system, the original five parameters can
be represented by three effective dimensionless combinations,
namely $\beta$, $\delta$, and $\mu$. In particular, the normalization
of the cubic coefficient is not an additional biological assumption;
rather, it corresponds to a choice of scales for the state variables.

For notational simplicity, in what follows we relabel
\[
(X,Y,\beta,\delta,\mu)
\longmapsto
(x,y,b,d,I),
\]
so that the normalized system is written as
\begin{equation}
\boxed{
\begin{aligned}
\dot{x}&=y-x^3+b x^2+I,\\
\dot{y}&=-d x^2-y.
\end{aligned}}
\tag{7}
\end{equation}

% ------------------------------------------------------------
\subsection{Biological interpretation of the normalization}
% ------------------------------------------------------------

The normalization above should be understood as a
nondimensionalization rather than as the biological assertion that
the original cubic coefficient is physically equal to one. The
coefficient $a$ determines the scale of the cubic nonlinear
contribution in the fast voltage equation. By rescaling the state
variables, its numerical value can be absorbed into the definition
of the dimensionless voltage and recovery variables.

Similarly, the translation
\begin{equation}
Y=\sqrt{a}\,(y-c)
\end{equation}
centers the recovery variable around the constant offset $c$. This
does not imply that the intrinsic parameter $c$ is biologically
absent. Rather, its constant contribution is absorbed into the choice
of the origin of the $y$-coordinate. In the transformed first
equation, $c$ and the applied current $I$ occur through the effective
combination
\begin{equation}
\mu=\sqrt{a}(c+I).
\end{equation}
Hence, from the viewpoint of the phase portrait of the reduced
two-dimensional autonomous system, $c$ and $I$ do not provide two
independent directions in the equilibrium equation.

Nevertheless, we retain the conventional biological interpretation
of $I$ as the externally applied current. Therefore, rather than
regarding the reduction as an elimination of the biological meaning
of $c$, we retain a biologically meaningful baseline value for $c$
when relating the reduced system to the standard Hindmarsh--Rose
parameterization.

% ------------------------------------------------------------
\subsection{Choice of fixed parameter values}
% ------------------------------------------------------------

For the subsequent analytical and numerical investigation, we fix
\begin{equation}
\boxed{
a=1,\qquad c=1,
}
\tag{8}
\end{equation}
and retain
\begin{equation}
\boxed{
b,\qquad d,\qquad I
}
\end{equation}
as the remaining parameters.

This choice has two complementary justifications. First, the
normalization $a=1$ is mathematically natural because the coefficient
of the cubic term can be absorbed by a rescaling of the state
variables when $a>0$, as demonstrated above. Second, the choice
$c=1$ provides a conventional dimensionless baseline for the
recovery equation and allows the reduced system to remain close to
the standard nondimensional Hindmarsh--Rose formulation.

Consequently, the model considered in the remainder of the analysis
is
\begin{equation}
\boxed{
\begin{aligned}
\dot{x}&=y-x^3+b x^2+I,\\
\dot{y}&=1-dx^2-y,
\end{aligned}}
\tag{9}
\end{equation}
where
\begin{equation}
\boxed{
(b,d,I)\in\mathbb{R}^3
}
\end{equation}
are treated as the effective parameters.

The parameter $b$ controls the quadratic contribution to the fast
nonlinear voltage equation and consequently influences the geometry
of the cubic nullcline. The parameter $d$ determines the strength of
the quadratic contribution in the recovery equation and therefore
affects the geometry of the recovery nullcline. Finally, $I$ retains
the interpretation of an externally applied input current and is
therefore particularly natural as a principal bifurcation parameter.

% ------------------------------------------------------------
\subsection{Why this parameterization is suitable for the present
analysis}
% ------------------------------------------------------------

The reduction from five parameters to the three-parameter family
\eqref{9} is advantageous for the analysis of limit cycles and global
phase portraits. In particular, the equilibrium equation becomes
\begin{equation}
\boxed{
x^3+(d-b)x^2-(1+I)=0,
}
\tag{10}
\end{equation}
which is a one-variable cubic equation. Thus the number and location
of finite equilibria can be determined explicitly in terms of
$(b,d,I)$.

Moreover, the Jacobian matrix at an equilibrium
$E^*=(x^*,y^*)$ is
\begin{equation}
J(E^*)
=
\begin{pmatrix}
-3(x^*)^2+2bx^* & 1\\
-2dx^* & -1
\end{pmatrix}.
\tag{11}
\end{equation}
Consequently,
\begin{equation}
\boxed{
\operatorname{tr}J(E^*)
=
-3(x^*)^2+2bx^*-1,
}
\tag{12}
\end{equation}
and
\begin{equation}
\boxed{
\det J(E^*)
=
x^*\left[3x^*+2(d-b)\right].
}
\tag{13}
\end{equation}

The Hopf condition is therefore given by
\begin{equation}
3(x_H)^2-2b x_H+1=0,
\tag{14}
\end{equation}
and hence
\begin{equation}
\boxed{
x_H^\pm
=
\frac{b\pm\sqrt{b^2-3}}{3}.
}
\tag{15}
\end{equation}
Thus the existence of real Hopf candidates requires
\begin{equation}
\boxed{
b^2\geq 3.
}
\tag{16}
\end{equation}

The corresponding values of the input current are obtained from the
equilibrium equation:
\begin{equation}
\boxed{
I_H^\pm
=
(x_H^\pm)^3
+
(d-b)(x_H^\pm)^2
-1.
}
\tag{17}
\end{equation}

Therefore, the reduced three-parameter system retains sufficient
freedom to investigate equilibrium multiplicity, Hopf bifurcations,
periodic orbits, and the associated global phase portraits, while
avoiding the substantial algebraic redundancy of the original
five-parameter representation.

% ------------------------------------------------------------
\subsection{Two-parameter global study}
% ------------------------------------------------------------

For a detailed two-parameter bifurcation and global phase-plane
analysis, one may subsequently fix the intrinsic parameter $d$ at
the conventional baseline value
\begin{equation}
\boxed{
d=5,
}
\tag{18}
\end{equation}
and regard
\begin{equation}
\boxed{
(b,I)
}
\end{equation}
as the principal bifurcation parameters. The resulting family is
\begin{equation}
\boxed{
\begin{aligned}
\dot{x}&=y-x^3+b x^2+I,\\
\dot{y}&=1-5x^2-y.
\end{aligned}}
\tag{19}
\end{equation}

The value $d=5$ is not introduced as a mathematical normalization;
rather, it is adopted as a biologically motivated baseline choice
consistent with the commonly used dimensionless Hindmarsh--Rose
parameterization. In this two-parameter setting, $b$ controls the
shape of the fast nonlinear response, while $I$ serves as the
externally controllable input current.

Accordingly, the analysis may be organized in two stages. First, the
three-parameter family \eqref{9} is used for the general analytical
derivation of equilibria, local stability and bifurcation conditions.
Second, the biologically motivated choice $d=5$ is employed for the
detailed exploration of the $(b,I)$ parameter plane, including limit
cycles and global phase portraits.

gggggggggggg

\documentclass[11pt,a4paper]{article}
\usepackage[utf8]{inputenc}
\usepackage[T1]{fontenc}
\usepackage{amsmath, amssymb, amsthm}
\usepackage{geometry}
\geometry{margin=1in}

\title{Simplified 2D Hindmarsh--Rose Model: \\ Finite Singular Points and Bifurcation Analysis}
\author{}
\date{}

\begin{document}

\maketitle

\section{The Simplified Model}

The general 2D Hindmarsh--Rose-type system is given by:
\begin{equation}
\begin{aligned}
\dot{x} &= y - ax^3 + bx^2 + I, \\
\dot{y} &= c - dx^2 - y,
\end{aligned}
\end{equation}
where \(x\) represents the membrane potential, \(y\) is the recovery variable, and \(I\) is the external input current. The parameters \(a, b, c, d\) control the intrinsic nonlinearities.

To obtain a mathematically tractable model for global phase-portrait analysis, we fix the parameters as follows:
\begin{itemize}
    \item \(a = 1\)
    \item \(c = 1\)
    \item \(b = d = k\) (where \(k\) is a fixed positive constant, biologically motivated by the standard Hindmarsh--Rose baseline \(k=5\))
\end{itemize}

The system reduces to:
\begin{equation}
\begin{aligned}
\dot{x} &= y - x^3 + kx^2 + I, \\
\dot{y} &= 1 - kx^2 - y.
\end{aligned} \tag{S}
\label{eq:simplified}
\end{equation}
Here, \(I\) is treated as the primary bifurcation parameter.

\section{Finite Singular Points}

\subsection{Solving for the Equilibria}
Set \(\dot{x} = 0\) and \(\dot{y} = 0\).
From \(\dot{y} = 0\), we obtain:
\begin{equation}
y = 1 - kx^2
\end{equation}
Substitute this into \(\dot{x} = 0\):
\begin{equation}
(1 - kx^2) - x^3 + kx^2 + I = 0
\end{equation}
The \(kx^2\) terms cancel out perfectly, yielding:
\begin{equation}
x^3 = I + 1
\end{equation}
For \(I > -1\), there exists a \textbf{unique real equilibrium}:
\begin{equation}
x^* = \sqrt[3]{I+1}, \qquad y^* = 1 - k(I+1)^{2/3}.
\end{equation}
This avoids the intractable cubic radicals associated with the general model.

\subsection{The Jacobian Matrix}
Let \(P(x,y) = y - x^3 + kx^2 + I\) and \(Q(x,y) = 1 - kx^2 - y\). The Jacobian matrix is:
\begin{equation}
J(x,y) = \begin{pmatrix} 
\frac{\partial P}{\partial x} & \frac{\partial P}{\partial y} \\ 
\frac{\partial Q}{\partial x} & \frac{\partial Q}{\partial y} 
\end{pmatrix} 
= \begin{pmatrix} 
-3x^2 + 2kx & 1 \\ 
-2kx & -1 
\end{pmatrix}.
\end{equation}
Evaluated at the equilibrium \((x^*, y^*)\), and letting \(u = x^* = (I+1)^{1/3}\), we have:
\begin{equation}
J^* = \begin{pmatrix} 
-3u^2 + 2ku & 1 \\ 
-2ku & -1 
\end{pmatrix}.
\end{equation}

\subsection{Trace and Determinant}
The trace \(T\) and determinant \(D\) of \(J^*\) are:
\begin{align}
T &= \operatorname{tr}(J^*) = -3u^2 + 2ku - 1, \\
D &= \det(J^*) = (-3u^2 + 2ku)(-1) - (1)(-2ku) = 3u^2 - 2ku + 2ku = 3u^2.
\end{align}
Notice that \(D = 3u^2 > 0\) for all \(I > -1\). Therefore, the equilibrium is \textbf{never a saddle}. It is either a node or a focus, and its stability is determined entirely by the sign of the trace \(T\).

\subsection{Classification of the Singular Point}
\begin{itemize}
    \item If \(T^2 - 4D > 0\): Node (stable if \(T < 0\), unstable if \(T > 0\)).
    \item If \(T^2 - 4D < 0\): Focus (stable if \(T < 0\), unstable if \(T > 0\)).
    \item If \(T = 0\) and \(D > 0\): Center (linearization) or weak focus.
\end{itemize}

\section{Hopf Bifurcation Analysis}
A Hopf bifurcation occurs when \(T = 0\) and \(D > 0\).
\begin{equation}
T = -3u^2 + 2ku - 1 = 0 \implies 3u^2 - 2ku + 1 = 0.
\end{equation}
Solving for \(u\):
\begin{equation}
u = \frac{2k \pm \sqrt{4k^2 - 12}}{6} = \frac{k \pm \sqrt{k^2 - 3}}{3}.
\end{equation}
For real roots, we require \(k^2 \ge 3 \implies |k| \ge \sqrt{3}\).

Adopting the biologically standard \(k=5\), we obtain:
\begin{equation}
u = \frac{5 \pm \sqrt{25-3}}{3} = \frac{5 \pm \sqrt{22}}{3}.
\end{equation}
Since \(5 > \sqrt{22}\), both roots are positive. The corresponding critical current values are:
\begin{equation}
I_{\text{Hopf}} = u^3 - 1 = \left(\frac{5 \pm \sqrt{22}}{3}\right)^3 - 1.
\end{equation}
This yields two explicit Hopf bifurcation values for \(I\), allowing for a clean analytical study of the transition from a stable equilibrium to limit cycles (spiking behavior) and back.

\section*{Why This Case is Ideal for a New Paper}
\begin{enumerate}
    \item \textbf{Tractable Algebra:} The unique equilibrium \(x^* = \sqrt[3]{I+1}\) is explicit for \(I > -1\), eliminating the need to solve a general cubic.
    \item \textbf{Clean Trace/Determinant:} The determinant \(D = 3u^2\) is always positive, so no saddle points exist in the biological region. Stability is governed entirely by the trace \(T\).
    \item \textbf{Explicit Hopf Bifurcations:} Two explicit Hopf bifurcation points in terms of \(I\) provide a clear route to studying limit cycles.
    \item \textbf{Poincar\'e Compactification:} The highest-degree terms (\(-x^3\) and \(-kx^2\)) preserve the essential nonlinear structure, making the behavior at infinity tractable while retaining exact control over the finite singular point.
    \item \textbf{Biological Meaning:} Setting \(b=d\) is a mathematical normalization that preserves the fast-slow structure and nonlinear spiking behavior of the HR neuron.
\end{enumerate}



\section{The Simplified Model}
We study the system
\begin{equation}
\begin{aligned}
\dot{x} &= y - x^3 + kx^2 + I, \\
\dot{y} &= 1 - kx^2 - y,
\end{aligned} \tag{S}
\label{eq:model}
\end{equation}
where \(k\) is a fixed positive constant (e.g., \(k=5\)) and \(I\) is the primary bifurcation parameter. The polynomial degree is \(n=3\).

\section{Poincar\'e Compactification}

To study the behaviour at infinity, we use the Poincar\'e compactification. We consider the local charts \(U_1\) and \(U_2\).

\subsection{Local Chart \(U_1\)}
In the local chart \(U_1\), we set
\[
x = \frac{1}{u}, \qquad y = \frac{v}{u}.
\]
After the change of variables and a time rescaling \(d\tau = dt/u^{n-1} = dt/u^2\), system \eqref{eq:model} becomes
\begin{equation}
\begin{aligned}
\dot{u} &= -u^3 v + u - k u^2 - I u^4, \\
\dot{v} &= u^3 - k u - u^2 v - u^2 v^2 + v - k u v - I u^3 v.
\end{aligned} \tag{U1}
\label{eq:U1}
\end{equation}

The line at infinity in this chart is \(u=0\). Setting \(u=0\) in \eqref{eq:U1} yields
\[
\dot{u} = 0, \qquad \dot{v} = v.
\]
Thus the only infinite singular point in \(U_1\) is
\[
(u,v) = (0,0).
\]
The Jacobian matrix of \eqref{eq:U1} at \((0,0)\) is
\[
J_{U_1}(0,0) = \begin{pmatrix} 1 & 0 \\ -k & 1 \end{pmatrix}.
\]
Its eigenvalues are both \(1\). Therefore, the infinite singular point \((0,0)\) in \(U_1\) is a \textbf{hyperbolic unstable node} (a source) on the equator. This point corresponds to the direction of the \(x\)-axis at infinity.

\subsection{Local Chart \(U_2\)}
In the local chart \(U_2\), we set
\[
x = \frac{u}{v}, \qquad y = \frac{1}{v}.
\]
After the change of variables and a time rescaling \(d\tau = dt/v^{n-1} = dt/v^2\), system \eqref{eq:model} becomes
\begin{equation}
\begin{aligned}
\dot{u} &= -u^3 + k u^3 v + k u^2 v + u v^2 - u v^3 + v^2 + I v^3, \\
\dot{v} &= v^3(1 - v + k u^2).
\end{aligned} \tag{U2}
\label{eq:U2}
\end{equation}

The line at infinity in this chart is \(v=0\). Setting \(v=0\) in \eqref{eq:U2} yields
\[
\dot{u} = -u^3, \qquad \dot{v} = 0.
\]
Thus the only infinite singular point in \(U_2\) is
\[
(u,v) = (0,0).
\]
The Jacobian matrix of \eqref{eq:U2} at \((0,0)\) is identically zero:
\[
J_{U_2}(0,0) = \begin{pmatrix} 0 & 0 \\ 0 & 0 \end{pmatrix}.
\]
Hence \((0,0)\) is a \textbf{linearly zero singular point}. Its local phase portrait must be determined by the blow-up technique. The leading terms near the origin are
\[
\dot{u} \approx -u^3 + v^2, \qquad \dot{v} \approx v^3.
\]
This is a quasi-homogeneous singular point of type \((3,2)\). A standard blow-up analysis (e.g., using weights \(u = r^2 X\), \(v = r^3 Y\)) shows that the origin is formed by a combination of hyperbolic and parabolic sectors. In particular, the flow on the equator \(v=0\) satisfies \(\dot{u} = -u^3\), so \(u=0\) is a stable equilibrium along the equator. Off the equator, \(\dot{v} = v^3\), so the flow moves away from the equator. The detailed sector decomposition yields \textbf{two hyperbolic sectors and two parabolic sectors}, similar to the even-\(\gamma\) case in the Higgins--Selkov system. This point corresponds to the direction of the \(y\)-axis at infinity.

\section{Summary of Infinite Singular Points}

\begin{itemize}
    \item In chart \(U_1\): \((0,0)\) is a hyperbolic unstable node (source) on the equator.
    \item In chart \(U_2\): \((0,0)\) is a linearly zero singular point, which after blow-up consists of two hyperbolic and two parabolic sectors.
\end{itemize}

These two points are the only infinite singular points on the Poincar\'e sphere. The qualitative behaviour at infinity is independent of the parameter \(I\) for \(I>-1\), since the leading terms are governed by the fixed cubic and quadratic nonlinearities.














\end{document}